% Galois + Spectral + Operads (Ch 12-14)
% Lines 2038-2432 of main.tex
% This file is for agent editing — will be merged back

\chapter{Galois Connections for Abstract Specification Inference}
\label{ch:galois}

Abstract interpretation (Cousot \& Cousot, 1977) computes sound
approximations of program behavior using \emph{abstract domains}.
CTTT incorporates abstract interpretation via Galois connections
between the concrete presheaf and abstract presheaves, dramatically
increasing the power of automatic specification inference.

\section{Abstract Domains as Presheaves}

\begin{definition}[Abstract Presheaf]
An \emph{abstract presheaf} $\Sem_f^\# : \site^{\mathrm{op}} \to
\mathsf{Lattice}$ assigns to each type fiber $U$ an element of an
abstract domain $D$, approximating the concrete behavior
$\Sem_f(U)$.

Examples of abstract domains:
\begin{center}
\begin{tabular}{l l l}
  \textbf{Domain} & \textbf{Elements} & \textbf{What it captures} \\
  \hline
  Sign & $\{-, 0, +, \top\}$ & Sign of integer output \\
  Interval & $[l, u] \subseteq \Int$ & Range of output \\
  Parity & $\{\mathsf{even}, \mathsf{odd}, \top\}$ & Parity \\
  Length & $\Nat \cup \{\top\}$ & Output collection length \\
  Sortedness & $\{\mathsf{sorted}, \mathsf{unsorted}, \top\}$ &
    Whether output is sorted \\
  Permutation & $\{\mathsf{perm}(x), \mathsf{not\_perm}, \top\}$ &
    Whether output permutes input \\
  Monotonicity & $\{\uparrow, \downarrow, \mathsf{const}, \top\}$ &
    Whether output is monotone in input \\
\end{tabular}
\end{center}
\end{definition}

\begin{definition}[Galois Connection]
A \emph{Galois connection} between concrete presheaf $\Sem_f$ and
abstract presheaf $\Sem_f^\#$ is a pair of monotone maps:
\begin{align}
  \alpha &: \Sem_f \to \Sem_f^\# && \text{(abstraction)} \\
  \gamma &: \Sem_f^\# \to \Sem_f && \text{(concretization)}
\end{align}
satisfying $\mathrm{id} \leq \gamma \circ \alpha$ and
$\alpha \circ \gamma \leq \mathrm{id}$ (the Galois condition).
\end{definition}

\begin{theorem}[Abstract Spec Inference via Galois Connection]
\label{thm:abstract-spec}
The abstract presheaf $\Sem_f^\#$ inferred by abstract interpretation
is a \emph{sound over-approximation} of the specification:
\[
  f \models S \implies \alpha(\Sem_f) \leq \alpha(\Sem_S)
\]

Contrapositively: if $\alpha(\Sem_f) \not\leq \alpha(\Sem_S)$, then
$f$ violates $S$ (sound bug detection).
\end{theorem}

\section{How This Increases CTTT's Power}

The basic CTTT (Parts I--III) compares Z3 terms directly, which
fails when terms are structurally different.  Abstract interpretation
\emph{lifts} the comparison to a coarser domain where structurally
different terms may have the same abstract value.

\begin{example}[Sign-Domain Spec Inference]
Consider:
\begin{lstlisting}
def f(n):        # recursive factorial
    if n <= 1: return 1
    return n * f(n-1)

def g(n):        # iterative factorial
    r = 1
    for i in range(2, n+1): r *= i
    return r
\end{lstlisting}

In the Sign domain:
\begin{itemize}
  \item $\alpha(\Sem_f)(U_\mathsf{Int}^{+}) = +$ (positive for positive $n$).
  \item $\alpha(\Sem_g)(U_\mathsf{Int}^{+}) = +$ (same).
  \item The abstract presheaves agree $\implies$ no obstruction
    at the abstract level.
\end{itemize}

This doesn't prove equivalence (the Sign domain is too coarse), but
it \emph{eliminates impossible counterexamples}: Z3 can be told
``the output is positive,'' narrowing its search.
\end{example}

\begin{example}[Sortedness-Domain Spec Checking]
\begin{lstlisting}
@guarantee("returns a sorted list")
def f(xs):
    result = []
    for x in xs:
        bisect.insort(result, x)
    return result
\end{lstlisting}

The Sortedness abstract domain infers:
$\alpha(\Sem_f) = \mathsf{sorted}$ (because $\mathsf{bisect.insort}$
maintains sorted invariant --- this is a \emph{transfer function}
in the abstract domain).

The specification $S$ requires $\alpha(\Sem_S) = \mathsf{sorted}$.
Since $\alpha(\Sem_f) = \alpha(\Sem_S)$, the spec is satisfied at
the abstract level.  This is a \emph{proof} (not just a test) that
the function returns a sorted list.
\end{example}

\section{The Galois--\v{C}ech Interaction}

The abstract domains interact with the \v{C}ech complex:

\begin{theorem}[Abstract Cohomology]
\label{thm:abstract-cohomology}
There is a morphism of cochain complexes:
\[
\begin{tikzcd}
  C^0(\covers, \Iso) \arrow[r, "\delta^0"] \arrow[d, "\alpha"] &
  C^1(\covers, \Iso) \arrow[r, "\delta^1"] \arrow[d, "\alpha"] &
  C^2(\covers, \Iso) \arrow[d, "\alpha"] \\
  C^0(\covers, \Iso^\#) \arrow[r, "\delta^{0,\#}"] &
  C^1(\covers, \Iso^\#) \arrow[r, "\delta^{1,\#}"] &
  C^2(\covers, \Iso^\#)
\end{tikzcd}
\]
The abstraction $\alpha$ commutes with the coboundary operators.
Therefore:
\[
  \Hone(\covers, \Iso^\#) = 0 \implies \Hone(\covers, \Iso) = 0
\]
(if the abstract \v{C}ech complex has trivial $\Hone$, so does the
concrete one --- but not vice versa).
\end{theorem}

This gives a \emph{cheaper} sufficient condition for equivalence:
compute $\Hone$ in the abstract domain (faster, always terminates)
before resorting to the full Z3 check (slower, may timeout).


\chapter{The Descent Spectral Sequence}
\label{ch:spectral}

The spectral sequence is a systematic tool for computing cohomology
in stages, each refining the previous.  It vastly increases CTTT's
power by enabling \emph{incremental} verification: start with a
coarse check, refine only where needed.

\section{Filtrations of the Type Site}

\begin{definition}[Type Complexity Filtration]
The type site $\site$ has a natural filtration by complexity:
\begin{align}
  F^0 &= \{\mathsf{None}, \mathsf{Bot}\}
    && \text{(trivial types)} \\
  F^1 &= F^0 \cup \{\mathsf{Bool}\}
    && \text{(finite types)} \\
  F^2 &= F^1 \cup \{\mathsf{Int}\}
    && \text{(countable types)} \\
  F^3 &= F^2 \cup \{\mathsf{Str}\}
    && \text{(string type)} \\
  F^4 &= F^3 \cup \{\mathsf{Pair}\}
    && \text{(structural types)} \\
  F^5 &= F^4 \cup \{\mathsf{Ref}\}
    && \text{(reference types --- heap)}
\end{align}
Each level adds more complex types, requiring more sophisticated
reasoning.
\end{definition}

\begin{definition}[The Spectral Sequence]
The filtration induces a spectral sequence $\{E_r^{p,q}\}$ with:
\begin{align}
  E_0^{p,q} &= C^{p+q}(F^p / F^{p-1}, \Iso) \\
  E_1^{p,q} &= H^{p+q}(F^p / F^{p-1}, \Iso) \\
  E_\infty^{p,q} &\implies H^{p+q}(\site, \Iso)
\end{align}
\end{definition}

\section{The Verification Strategy}

The spectral sequence dictates an \emph{optimal verification strategy}:

\begin{enumerate}
  \item \textbf{$E_0$ page}: check equivalence on each filtration
    \emph{grade} independently.  For the None/Bot grade: trivially
    equivalent (both return the same constant on trivial inputs).
    For the Bool grade: finitely many inputs, decidable.
    For the Int grade: Z3 integer arithmetic, often decidable.

  \item \textbf{$E_1$ page}: compute $\Hone$ within each grade.
    If $\Hone(F^p / F^{p-1}) = 0$ for all $p$, the spectral
    sequence \emph{degenerates} and $\Hone(\site) = 0$ (equivalence).

  \item \textbf{$E_2$ page}: if $E_1$ has non-trivial differentials,
    compute $E_2 = H(E_1, d_1)$.  This captures \emph{interactions}
    between adjacent grades (e.g., an integer function whose behavior
    depends on a boolean flag).

  \item \textbf{Higher pages}: each $E_r$ refines the approximation.
    In practice, the spectral sequence degenerates at $E_1$ or $E_2$
    for most programs (because the type fibers are independent).
\end{enumerate}

\begin{theorem}[Spectral Sequence Degeneration for Duck-Typed Programs]
\label{thm:degeneration}
If duck type inference determines that each parameter has a
\emph{single} type (the most common case), the spectral sequence
degenerates at $E_1$:
\[
  E_1^{p,q} = E_\infty^{p,q} \quad \text{for all $p, q$}
\]
and $\Hone(\site, \Iso) = \bigoplus_p E_1^{p,0}$ (direct sum of
fiber-wise cohomologies).
\end{theorem}

\begin{proof}
When each parameter has a unique type, the type site has a single
object (one fiber combination).  There are no overlaps, no transition
functions, and the \v{C}ech complex is trivial: $C^0 = \Iso(U)$,
$C^1 = C^2 = 0$.  Therefore $\Hone = 0$ (no obstructions to check)
and the local verdict IS the global verdict.  The spectral sequence
has a single non-trivial entry $E_1^{0,0} = H^0(U, \Iso)$ and
degenerates immediately.
\end{proof}

This theorem explains why our current implementation achieves good
results on programs with clear types: the spectral sequence
degenerates, and a single Z3 query per fiber suffices.

\section{Non-Degenerate Cases: Type Polymorphism}

When parameters are polymorphic (duck-typed with multiple possible
types), the spectral sequence does NOT degenerate, and the $E_2$
page reveals interactions:

\begin{example}[Type-Dependent Behavior]
\begin{lstlisting}
def f(x):
    if isinstance(x, int):
        return x * 2
    elif isinstance(x, str):
        return x + x
    else:
        return [x, x]
\end{lstlisting}

The type site has three fibers: Int, Str, Other.  On each fiber,
$f$ has different behavior.  The $E_1$ page checks each fiber
independently.  The $E_2$ page checks compatibility across fibers:
does the ``doubling'' specification ($S(x, y) = (y = x + x)$)
make sense across type boundaries?

For integers, $x + x = 2x$ (integer doubling).
For strings, $x + x = xx$ (string repetition).
These are \emph{different operations} ($\mathsf{ADD}$ vs.\
$\mathsf{Concat}$) that happen to share the $+$ syntax.

The $E_2$ differential detects this: the integer and string fibers
use different path constructors ($\mathsf{add\_comm}$ vs.\ string
concat), and the transition function on the Int$\cap$Str overlap
is non-trivial.  This is a genuinely \emph{higher} cohomological
phenomenon --- invisible to $E_1$.
\end{example}


\chapter{Operad Theory for Higher-Order Programs}
\label{ch:operads}

Python programs frequently use \emph{higher-order functions}:
functions that take other functions as arguments (\texttt{map},
\texttt{filter}, \texttt{reduce}, \texttt{sorted} with key, etc.).
The CTTT framework handles these via \emph{operad theory}, which
provides the algebraic structure for composing operations.

\section{The Operad of Python Operations}

\begin{definition}[The Python Operations Operad]
The \emph{Python operations operad} $\mathcal{O}_\mathsf{Py}$ has:
\begin{itemize}
  \item \textbf{Colors} (sorts): $\PyObj$, $\Bool$, $\Int$,
    $\mathsf{Str}$, $\mathsf{Seq}$, $\mathsf{Map}$.
  \item \textbf{Operations}: for each Python builtin/method, an
    operation $o \in \mathcal{O}(c_1, \ldots, c_n; c)$ mapping
    $n$ inputs of colors $c_1, \ldots, c_n$ to an output of color $c$.
  \item \textbf{Composition}: operadic composition $\circ_i$
    (substituting an operation into the $i$-th input of another).
  \item \textbf{Equivariance}: permutation of inputs (for
    commutative operations).
\end{itemize}
\end{definition}

\begin{example}[Operations in the Operad]
\begin{align}
  + &\in \mathcal{O}(\Int, \Int; \Int) \\
  \mathsf{sorted} &\in \mathcal{O}(\mathsf{Seq}; \mathsf{Seq}) \\
  \mathsf{map} &\in \mathcal{O}((\PyObj \to \PyObj), \mathsf{Seq};
    \mathsf{Seq}) \\
  \mathsf{reduce} &\in \mathcal{O}((\PyObj \times \PyObj \to \PyObj),
    \PyObj, \mathsf{Seq}; \PyObj) \\
  \mathsf{filter} &\in \mathcal{O}((\PyObj \to \Bool), \mathsf{Seq};
    \mathsf{Seq})
\end{align}
\end{example}

\section{Operadic Path Constructors}

The operad structure provides new path constructors beyond the
basic algebraic axioms:

\begin{theorem}[Map--Fold Fusion]
\label{thm:map-fold}
For a monoid $(M, \oplus, e)$ and a function $h : A \to M$:
\[
  \mathsf{reduce}(\oplus, e, \mathsf{map}(h, \mathsf{xs}))
  \equiv
  \mathsf{reduce}(\lambda \mathsf{acc}, x.\; \mathsf{acc} \oplus h(x),
    e, \mathsf{xs})
\]
This is the \emph{map--fold fusion law}: mapping then folding is
equivalent to folding with a composed operation.
\end{theorem}

\begin{proof}
By induction on $\mathsf{xs}$ (Nat-recursor uniqueness applied to
the list length):
\begin{itemize}
  \item Base: both sides return $e$ on the empty list.
  \item Step: the left side maps $h$ to the head, then folds; the
    right side applies $\oplus$ with $h(\mathsf{head})$ directly.
    Both produce $\mathsf{acc} \oplus h(\mathsf{head})$.
\end{itemize}
\end{proof}

\begin{theorem}[Filter--Map Commutation]
\label{thm:filter-map}
If the predicate $p$ and the transformation $h$ are independent
($p(h(x)) = p(x)$ for all $x$):
\[
  \mathsf{filter}(p, \mathsf{map}(h, \mathsf{xs}))
  \equiv
  \mathsf{map}(h, \mathsf{filter}(p, \mathsf{xs}))
\]
\end{theorem}

\begin{theorem}[Sorted--Map Interaction]
\label{thm:sorted-map}
If $h$ is \emph{monotone} ($x \leq y \implies h(x) \leq h(y)$):
\[
  \mathsf{sorted}(\mathsf{map}(h, \mathsf{xs}))
  \equiv
  \mathsf{map}(h, \mathsf{sorted}(\mathsf{xs}))
\]
Sorting then mapping a monotone function is the same as mapping
then sorting.
\end{theorem}

These operadic laws generate new paths in the HIT $\mathsf{PyComp}$,
enabling equivalence proofs for programs that compose higher-order
operations in different ways.

\section{Application: Higher-Order Equivalence}

\begin{example}[Map+Filter vs. Comprehension with Condition]
\begin{lstlisting}
# Program A: map then filter
def f(xs):
    return list(filter(lambda x: x > 0, map(lambda x: x*2, xs)))

# Program B: comprehension with condition
def g(xs):
    return [x*2 for x in xs if x*2 > 0]
\end{lstlisting}

The operadic path:
\begin{enumerate}
  \item $\mathsf{filter}(p, \mathsf{map}(h, \mathsf{xs}))$
    (program A).
  \item By filter--map commutation (since $p(h(x)) = (h(x) > 0)$
    and $p(x) = (x > 0)$ --- but $h(x) = 2x$, so
    $p(h(x)) = (2x > 0) \neq (x > 0)$ in general!  The commutation
    does NOT apply here).
  \item Instead, use the comprehension--loop equivalence (D4):
    the comprehension $[h(x) \text{ for } x \text{ if } p(h(x))]$
    is equivalent to $\mathsf{filter}(p, \mathsf{map}(h, \mathsf{xs}))$.
  \item The path is $\mathsf{D4}$ (comprehension fusion).
\end{enumerate}
\end{example}


